\section{Discussion}
\label{sec:discussion}

\paragraph{Three failures with three different remedies.} The audit's value is that it separates
things that are usually lumped together under "reproducibility".

\emph{A number that reproduces on some datasets and misses on others} is harder to diagnose than one
that never reproduces, because the environment is demonstrably working: the same SpikingGCN script
reaches the published CiteSeer value and misses the Cora value by thirty points. The remedy is a
per-dataset reporting convention and a released configuration, not a global environment fix.

\emph{A number that reproduces on a split nobody else uses} is more insidious, because nothing is
broken. The code is correct, the script is honest, the number is real; it simply measures a different
task from the one the reader will compare it with, and the difference is ten points. The remedy is not
more code but a reporting convention: a table entry should carry its split, its number of runs, and
whether the split is the standard one, in the table rather than in a footnote.

\emph{An unlabelled split choice} is the most damaging, because it corrupts a comparison table
without anyone making a mistake. SpikingGCN's own paper prints $80.7$ on Split~I and $88.7$ on
Split~II; the later tables cite $81.5$ and $87.7$ without saying which. Both numbers are correct.
The remedy is a reporting convention: a baseline row should name its split, and re-running a baseline
should be the default rather than the exception.

\paragraph{What a reader should do.} Practically: do not assemble a baseline table from published
numbers in this area. Pick the method you intend to compare against, run its code on your split, and
report that. Each pipeline takes one working session to re-run once its environment is pinned,
which is what makes the recommendation practical.

\paragraph{Why the field is unusually exposed.} Two structural facts make spiking GNNs worse than the
average subfield. The first is that the spiking encoder is a stochastic sampler reading its input as a
probability, so a preprocessing choice that is cosmetic in a conventional GNN --- row-normalising the
features --- silently disables learning rather than degrading it. The second is that a spiking
network's accuracy depends on the number of timesteps, a hyper-parameter that is reported as part of
the method's efficiency claim rather than as part of its accuracy claim, so two numbers can differ by
ten points without either paper being wrong about what it measured.

\paragraph{Division of evidence with the companion paper.} The conversion study \citep{companion}
uses two numbers from this audit --- $74.0$ and $86.6$ on the public split --- as an external anchor
for its cost comparison. It does not report the audit protocol, the reproduction grid or the
split-mixing analysis, and none of its claims depends on them. This paper reports those results and
does not use the conversion method's fidelity, cost or theory results. The overlap between the two
manuscripts is therefore two measured numbers, explicitly cross-cited, rather than a shared result
set.

\paragraph{Threats to this audit.} Ours, not the audited papers'. \emph{Environment drift}: where a
stated dependency no longer exists for our platform we installed the nearest available version, and
recorded it; a reader with the original stack may get different numbers, and if they do, that is
itself the point. \emph{Under-powered cells}: four of the seven rows are single-run, and we draw no
conclusion from a difference smaller than the spread the multi-seed rows show. \emph{Cora-centred}:
four of seven rows are Cora, chosen because it is the dataset all three papers report, not because it
is representative; a dataset where all three methods agree might behave differently. \emph{Our own
code is in the comparison}: the $80.5$ in \secref{sec:results} comes from a companion paper, and a
reader should discount it accordingly; it is not load-bearing for any claim here.

\paragraph{Evaluative claims and their scope.} The audit supports one claim: a published accuracy in
this area is not transportable to a different split or codebase without being re-run. It does not
support a ranking of the three methods, a judgement on the correctness of any audited paper, or a
claim about spiking GNNs beyond the citation graphs measured here. That claim rests on three
assumptions --- that the released repository is the artefact the paper describes, that the pinned
environment reproduces the authors' execution conditions, and that the numbers in a comparison table
belong to the methods they name. Each is stated because each is checkable, and \secref{sec:mixed} is
what happens when the third fails.

\paragraph{Broader impact and intended use.} An executable audit is easy to misread as a quality
ranking, so every reproduced value here carries its split, its run count and the configuration that
produced it, and no audited paper is described as wrong. The constructive use is prospective: a new
method can be compared against a number produced on its own split by running the released code,
instead of against a figure copied from a table. This study introduces no new dataset --- all three
graphs are the public Planetoid splits used by the audited papers --- and the released artefact is a
protocol with its logs rather than a data resource.

\paragraph{What would falsify us.} Each of the three findings is a claim that a reader can check
against shipped artefacts. If SpikingGCN's code produces $80.7$ on Cora for you, our environment is
wrong and we would like to know. If SpikeGraphormer's public-split configuration produces $84.8$ for
you, we have mis-set a flag. If a later table labels which split it is quoting, the split-mixing
finding no longer applies to that table.
